% !TEX root = ./ECEN-635 - Homework 1 - Brody Jordan.tex

\documentclass[11pt]{article}

\usepackage[
    assignmentDueDate={September 12, 2025 at 4:30 pm}
]{C:/Users/bljor/Sync/Courses/assignment}

\usepackage{bm}
\usepackage{siunitx}
\usepackage{mathrsfs}
\usepackage{cases}

\begin{document}
\makeAssignmentTitle

% --- Problem 1 ---
\begin{homeworkProblem}

Convert the following time domain expressions for the electric and magnetic field into frequency domain expressions. 
Show each mathematical step in your calculations. A miraculous answer without the complete calculation will be judged harshly.

\begin{enumerate}[label=(\alph*), topsep=5pt, itemsep=0pt]
    \item $\mathbf{E}(t)=4 \cos (\omega t-28 z) \hat{\mathbf{y}} \;\; \unit{V/m}$
    \item $\mathbf{E}(t)=9 \sin (\omega t-23 z) \hat{\mathbf{x}} \;\; \unit{V/m}$
    \item $\mathbf{H}(t)=7 \cos (\omega t-23 z) \hat{\mathbf{y}} \;\; \unit{mA/m}$ \qquad \textit{changed polarization from x to y}
    \item $\mathbf{E}(t)=4.6 e^{-48 \pi z} \cos (\omega t-139 z+\pi / 4) \hat{\mathbf{y}} \;\; \unit{V/m}$
\end{enumerate} \vspace{1\baselineskip}

\solution

\textbf{Part (a)}
\begin{align*}
    \cos (\omega t - 28 z) &= \Re \{ e^{j(\omega t - 28 z)}\} \\
    &= \Re \{ e^{j\omega t} e^{-j 28 z}\} \\
    &\Rightarrow \quad \mathbf{E}(\omega) = 4 e^{-j 28 z} \hat{\mathbf{y}} \;\; \unit{V/m}
\end{align*}

\textbf{Part (b)}
\begin{align*}
    9 \sin (\omega t - 23 z) &= -9 \cos (\omega t - 23 z + \pi/2) \\
    &= \Re \{ -9 e^{j(\omega t - 23 z + \pi/2)}\} \\
    &\Rightarrow \quad \mathbf{E}(\omega) = 9 e^{j \pi/2} e^{-j 23 z} \hat{\mathbf{x}} \;\; \unit{V/m}
\end{align*}

\textbf{Part (c)}
\begin{align*}
    7 \cos (\omega t - 23 z) &= \Re \{ 7 e^{j(\omega t - 23 z)}\} \\
    &= \Re \{ 7 e^{j\omega t} e^{-j 23 z}\} \\
    &\Rightarrow \quad \mathbf{H}(\omega) = 7 e^{-j 23 z} \hat{\mathbf{y}} \;\; \unit{mA/m}
\end{align*}

\textbf{Part (d)}
\begin{align*}
    4.6 e^{-48 \pi z} \cos (\omega t - 139 z + \pi/4) &= 4.6 e^{-48 \pi z} \Re \{ e^{j(\omega t - 139 z + \pi/4)}\} \\
    &= 4.6 e^{-48 \pi z} \Re \{ e^{j\omega t} e^{-j 139 z} e^{j \pi/4}\} \\
    &\Rightarrow \quad \mathbf{E}(\omega) = 4.6 e^{-(48\pi +j 139) z} e^{j \pi/4} \hat{\mathbf{y}} \;\; \unit{V/m}
\end{align*}

\end{homeworkProblem}

\pagebreak

% --- Problem 2 ---
\begin{homeworkProblem}

Convert the following frequency domain expressions to the time domain. Show each mathematical step in your calculation.

\begin{enumerate}[label=(\alph*), topsep=5pt, itemsep=0pt]
    \item $\mathbf{E} = 15 e^{-j 23 z} \hat{\mathbf{x}} \;\; \unit{V/m}$
    \item $\mathbf{E} = 32 e^{-(52 + j 135) z} \hat{\mathbf{y}} \;\; \unit{V/m}$
    \item $\mathbf{H} = [(2 + j2) \hat{\mathbf{x}} + (1 - j 3) \hat{\mathbf{y}}] e^{-j 245 z} \;\; \unit{A/m}$
    \item $\rho = 3 j e^{j\pi/2} \;\; \unit{C/m^3}$ \\
    Be sure to place each of the vector components in simplest magnitude and phase form before
    converting to the time domain.
\end{enumerate} \vspace{1\baselineskip}

\solution

\textbf{Part (a)}
\begin{align*}
    \mathbf{E} &= 15 e^{-j 23 z} \hat{\mathbf{x}} \;\; \unit{V/m} \\
    &= \Re \{15 e^{j(\omega t - 23 z)}\} \hat{\mathbf{x}} \;\; \unit{V/m} \\
    &\Rightarrow \quad \mathbf{E}(t) = 15 \cos (\omega t - 23 z) \hat{\mathbf{x}} \;\; \unit{V/m}
\end{align*}

\textbf{Part (b)}
\begin{align*}
    \mathbf{E} &= 32 e^{-(52 + j 135) z
} \hat{\mathbf{y}} \;\; \unit{V/m} \\
    &= 32 e^{-52 z} e^{-j 135 z} \hat{\mathbf{y}} \;\; \unit{V/m} \\
    &= \Re \{32 e^{-52 z} e^{j(\omega t - 135 z)}\} \hat{\mathbf{y}} \;\; \unit{V/m} \\
    &\Rightarrow \quad \mathbf{E}(t) = 32 e^{-  52 z} \cos (\omega t - 135 z) \hat{\mathbf{y}} \;\; \unit{V/m}
\end{align*}

\textbf{Part (c)}
\begin{align*}
    \mathbf{H} &= [(2 + j2) \hat{\mathbf{x}} + (1 - j 3) \hat{\mathbf{y}}] e^{-j 245 z} \;\; \unit{A/m} \\
    &= [2\sqrt{2} e^{j \pi/4} \hat{\mathbf{x}} + \sqrt{10} e^{-j \arctan(3)} \hat{\mathbf{y}}] e^{-j 245 z} \;\; \unit{A/m} \\
    &= \Re \{2\sqrt{2} e^{j(\omega t - 245 z + \pi/4)} \hat{\mathbf{x}} + \sqrt{10} e^{j(\omega t - 245 z - \arctan(3))} \hat{\mathbf{y}}\} \;\; \unit{A/m} \\
    &\Rightarrow \quad \mathbf{H}(t) = 2\sqrt{2} \cos (\omega t - 245 z + \pi/4) \hat{\mathbf{x}} + \sqrt{10} \cos (\omega t - 245 z - \arctan(3)) \hat{\mathbf{y}} \;\; \unit{A/m}
\end{align*}

\textbf{Part (d)}
\begin{align*}
    \rho &= 3 j e^{j\pi/2} \;\; \unit{C/m^3} \\
    &= \Re \{3 e^{j(\omega t + \pi/2 + \pi/2)}\} \;\; \unit{C/m^3} \quad\quad \text{since } j=e^{j\pi/2} \\
    &\Rightarrow \quad \rho(t) = 3 \cos (\omega t + \pi) \;\; \unit{C/m^3}
\end{align*}

\end{homeworkProblem}

\pagebreak

% --- Problem 3 ---
\begin{homeworkProblem}

Assuming a homogeneous region where there is no current, charge, or loss:
\begin{enumerate}[label=(\alph*), topsep=5pt, itemsep=0pt]
    \item Use Maxwell's equations to convert the frequency domain electric field expression in 2(b) to the corresponding
    magnetic field. Show all mathematical steps.
    \item Use Maxwell’s equations to convert the frequency domain magnetic field in 2(c) into a frequency
    domain electric field into a frequency domain magnetic field. Show all mathematical steps.
\end{enumerate} \vspace{1\baselineskip}

\solution

\textbf{Part (a)} \vspace{1\baselineskip}

We have Maxwell's equations which provides a relationship between the electric and magnetic fields, and we have an electric field from 2(b): 
\[\begin{cases}
    \nabla \times \mathbf{E} = -j \omega \mathbf{H} \\
    \mathbf{E} = 32 e^{-(52 + j 135) z} \hat{\mathbf{y}}
\end{cases}\]
Since this electric field has only a y-component, the curl becomes
\[
    \nabla\times\mathbf{E} = \left(-\frac{\partial}{\partial z} E_y\right)\hat{\mathbf{x}} + \left(\frac{\partial}{\partial x} E_y\right) \hat{\mathbf{z}}
\]
The partial derivative of $E_y$ with respect to $x$ is zero, so our curl simplifies to
\begin{align*}
    \nabla\times\mathbf{E} &= -\frac{\partial}{\partial z} \left(32 e^{-(52 + j 135) z}\right) \hat{\mathbf{x}} + 0 \hat{\mathbf{z}} \\
    &= -32 (-(52 + j 135)) e^{-(52 + j 135) z} \hat{\mathbf{x}} \\
    &= 32 (52 + j 135) e^{-(52 + j 135) z} \hat{\mathbf{x}}
\end{align*}
Now, substituting this into Maxwell's equation, we can find the magnetic field
\begin{align*}
    32 (52 + j 135) e^{-(52 + j 135) z} \hat{\mathbf{x}} = -j \omega \mathbf{H} \\
    \implies \mathbf{H} = -\frac{32 (52 + j 135)}{j \omega} e^{-(52 + j 135) z} \hat{\mathbf{x}}
\end{align*}
\vspace{1\baselineskip}

\textbf{\textit{Part (b) continued on next page...}}

\pagebreak

\textbf{Part (b)} \vspace{1\baselineskip}

Now we have a magnetic field from 2(c) and Maxwell's equations:
\[\begin{cases}
    \nabla \times \mathbf{H} = j \omega \epsilon \mathbf{E} \\
    \mathbf{H} = [(2 + j2) \hat{\mathbf{x}} + (1 - j 3) \hat{\mathbf{y}}] e^{-j 245 z}
\end{cases}\]
Calculating the curl of $\mathbf{H}$, we have
\begingroup
\addtolength{\jot}{1em}
\begin{align*}
    \nabla\times\mathbf{H} &= \left(-\frac{\partial}{\partial z} H_y\right) \mathbf{\hat{x}} + \left(\frac{\partial}{\partial z} H_x\right) \mathbf{\hat{y}} + \left(\frac{\partial}{\partial x} H_y - \frac{\partial}{\partial y} H_x\right) \mathbf{\hat{z}} \\
    &= \left(-\frac{\partial}{\partial z} H_y\right) \mathbf{\hat{x}} + \left(\frac{\partial}{\partial z} H_x\right) \mathbf{\hat{y}} \\
    &= \left(-\frac{\partial}{\partial z} \left((1 - j 3) e^{-j 245 z}\right)\right) \mathbf{\hat{x}} + \left(\frac{\partial}{\partial z} \left((2 + j2) e^{-j 245 z}\right)\right) \mathbf{\hat{y}} \\
    &= \left(- (1 - j 3)(-j 245) e^{-j 245 z}\right) \mathbf{\hat{x}} + \left((2 + j2)(-j 245) e^{-j 245 z}\right) \mathbf{\hat{y}} \\
    &= \left(735 + j245\right) e^{-j 245 z} \mathbf{\hat{x}} + \left(490 - j 490\right) e^{-j 245 z} \mathbf{\hat{y}}
\end{align*}
\endgroup
Now relating it to the electric field using Maxwell's equations
\[
    \left(735 + j245\right) e^{-j 245 z} \mathbf{\hat{x}} + \left(490 - j 490\right) e^{-j 245 z} \mathbf{\hat{y}} = j \omega \epsilon \mathbf{E} \\
\]
\[
    \implies \mathbf{E} = \frac{1}{j \omega \epsilon} \left[ \left(735 + j245\right) e^{-j 245 z} \mathbf{\hat{x}} + \left(490 - j 490\right) e^{-j 245 z} \mathbf{\hat{y}} \right]
\]
\end{homeworkProblem}

\pagebreak

% --- Problem 4 ---
\begin{homeworkProblem}

Use 1(b) and the corresponding electric field found in problem 1(c) to calculate the time average power density. 
Show all mathematical steps. \vspace{1\baselineskip}

\solution

We have the electric and magnetic fields from problems 1(b) and 1(c) and a relationship between them:
\[\begin{cases}
    \mathbf{E}(t) = 9 \sin (\omega t - 23 z) \hat{\mathbf{x}} \;\; \unit{V/m} &\quad\quad \text{from 1(b)}  \\
    \mathbf{H}(t) = 7 \cos (\omega t - 23 z) \hat{\mathbf{y}} \;\; \unit{mA/m} &\quad\quad \text{from 1(c)} \\
    \bm{\mathscr{S}}_{av} = \mathbf{S} = \frac{1}{2} \Re \{\mathbf{E} \times \mathbf{H}^*\}
\end{cases}\]
In the frequency domain and in base SI units, we have
\[\begin{cases}
    \mathbf{E}(\omega) = 9 e^{-j \pi/2} e^{-j 23 z} \hat{\mathbf{x}} &\quad\quad \text{from 1(b)}  \\
    \mathbf{H}(\omega) = 7\cdot10^{-3} e^{-j 23 z} \hat{\mathbf{y}} &\quad\quad \text{from 1(c)}
\end{cases}\]
Now we can compute $\mathbf{E}\times \mathbf{H}^*$. Their cross will only have a $z$-component since $\hat{\mathbf{x}} \times \hat{\mathbf{y}} = \hat{\mathbf{z}}$.
\begin{align*}
    \mathbf{E} \times \mathbf{H}^* &= (9 e^{-j \pi/2} e^{-j 23 z} \hat{\mathbf{x}}) \times (7\cdot10^{-3} e^{j 23 z} \hat{\mathbf{y}}) \\
    &= 63\cdot10^{-3} e^{-j \pi/2} (\hat{\mathbf{x}} \times \hat{\mathbf{y}}) \\
    &= 63\cdot10^{-3} e^{-j \pi/2} \hat{\mathbf{z}} \\
    &= -j 63\cdot10^{-3} \hat{\mathbf{z}}
\end{align*}
So the time average power density is
\begin{align*}
    \mathbf{S} = \frac{1}{2} \Re \{\mathbf{E} \times \mathbf{H}^*\} &= \frac{1}{2} \Re \{-j 63\cdot10^{-3} \hat{\mathbf{z}}\} \\
    &= \frac{1}{2} (0) \hat{\mathbf{z}} \\
    &= 0 \hat{\mathbf{z}} \;\; \unit{W/m^2}
\end{align*}
\end{homeworkProblem}

\pagebreak

% --- Problem 5 ---
\begin{homeworkProblem}

Starting with the frequency domain form of Maxwell's equations, derive the magnetic field vector Helmholtz equation. Number
your Maxwell equations and at every step write briefly which Maxwell equation and/or which identity you are using
from the Vectory Identities section of Appendix II.3.1 (pg. 958). Assume there is no electric or magnetic current and
no electric or magnetic charge. \vspace{1\baselineskip}

\solution

Here are Maxwell's equations with no charges
\begin{numcases}{}
    \nabla\times\mathbf{E} = -j \omega \mu \mathbf{H} &\quad\quad \\
    \nabla\times\mathbf{H} = j \omega \epsilon \mathbf{E} &\quad\quad \\
    \nabla \cdot \mathbf{E} = 0 &\quad\quad \\
    \nabla \cdot \mathbf{H} = 0 &\quad\quad
\end{numcases}

Start by taking the curl of equation (2)
\begin{equation}
\nabla \times (\nabla \times \mathbf{H}) = \nabla \times (j \omega \epsilon \mathbf{E}) \quad\quad
\end{equation}
Using vector identity (Appendix II-51) on the left hand side of (5)
\begin{equation}
\nabla \times (\nabla \times \mathbf{H}) = \nabla(\nabla \cdot \mathbf{H}) - \nabla^2 \mathbf{H} \quad\quad
\end{equation}
We substitute (6) into (5) to get
\begin{equation}
\nabla(\nabla \cdot \mathbf{H}) - \nabla^2 \mathbf{H} = \nabla \times (j \omega \epsilon \mathbf{E})
\end{equation}
Now substituting equation (1) and (4) into (7)
\begin{equation}
- \nabla^2 \mathbf{H} = j \omega \epsilon (-j \omega \mu \mathbf{H})
\end{equation}
Simplifying (8) by grouping the constants into $k^2$ gives us 
\begin{equation}
\nabla^2 \mathbf{H} + \omega^2 \mu \epsilon \mathbf{H} = 0 \quad\quad \Rightarrow \quad\quad \nabla^2 \mathbf{H} + k^2 \mathbf{H} = 0
\end{equation}
Which is the magnetic field vector Helmholtz equation. :)

\end{homeworkProblem}

\end{document}